\section{Strategy 2: Exact Demand Propagation}
\label{sec:strategy2-reader}

This section presents the reusable mechanism without repeating separate CE1
and CE2 coordinate systems.  Section~\ref{sec:signed-ce12-reductions}
develops the signed center model and the local propagation interface; the
exact admissible-set catalogue and the remaining scalar certificates are in
Appendix~\ref{sec:exact-local-demand-calculus}.

We now build on the maps $\mathcal A$, $B_c$, $F_c$, and $G_c$ defined in
that section.  In particular, every composition below is interpreted through
the actual-row implication \eqref{eq:reader-actual-propagation}.

\subsection{The universal selected-$T_+$ curve}

The exact capped map has four genuine labels: Full, $L$, $T_-$, and selected
$T_+$.  Its only nonlinear transition is independent of the row deficit.

\begin{proposition}[Universal selected-$T_+$ curve]
\label{prop:reader-universal-tplus}
Fix $0<d<1/2$, put $c=1-d$, and suppose an input $p$ belongs to a selected
$T_+$ arc of $G_c$.  Write
\[
 q=G_c(p),
 \qquad
 \nu=q-p,
 \qquad
 x=\frac{q-d}{1-d}.
\]
Then
\begin{equation}
 \nu=\psi(x):=1-\sqrt{1-x+x^2},
 \label{eq:reader-psi}
\end{equation}
and
\begin{equation}
 q=d+(1-d)x,
 \qquad
 p=d+(1-d)x-\psi(x).
 \label{eq:reader-tplus-normal}
\end{equation}
Consequently $p\mapsto G_{1-d}(p)$ is increasing and strictly concave on
every selected $T_+$ arc.
\end{proposition}

\begin{proof}
The selected quadratic is
\[
 (1-q)(q-d)=(1-d)^2\nu(2-\nu).
\]
Substituting $q=d+(1-d)x$ gives
\[
 x(1-x)=\nu(2-\nu).
\]
The selected component uses the smaller root, which is
\eqref{eq:reader-psi}.  If $r(x)=1-x+x^2$, then
\[
 \psi'(x)=\frac{1-2x}{2\sqrt{r(x)}},
 \qquad
 \psi''(x)=-\frac{3}{4r(x)^{3/2}}<0.
\]
Thus
\[
 p_d(x)=d+(1-d)x-\psi(x)
\]
is increasing and strictly convex.  Its inverse is increasing and strictly
concave, while $q=d+(1-d)x$ is affine in that inverse.
\end{proof}

With $e(d)$ defined by \eqref{eq:signed-low-root-definition}, the selected
high-radial arc has endpoints
$(d,d)$ and $(e(d),d+e(d))$, so concavity gives
\begin{equation}
 G_{1-d}(p)
 \ge p+\frac{d}{e(d)-d}(p-d).
 \label{eq:reader-high-chord}
\end{equation}
In the low-radial range, if $t=h(1-d)$ is the order-transition value, the
endpoints are $(d,d)$ and $(t,1-t)$, and
\begin{equation}
 G_{1-d}(p)
 \ge p+\frac{1-2t}{t-d}(p-d).
 \label{eq:reader-low-chord}
\end{equation}

For algebraic calculations put $z=\psi(x)/x$.  Then
\begin{equation}
 x=\frac{1-2z}{1-z^2},
 \qquad
 \psi(x)=\frac{z(1-2z)}{1-z^2}.
 \label{eq:reader-tplus-rational}
\end{equation}
In the historical notation $x=1-\beta$ and
$m_\beta=\sqrt{\beta^2-\beta+1}$,
\begin{equation}
 \beta=\frac{z(2-z)}{1-z^2},
 \qquad
 m_\beta=\frac{1-z+z^2}{1-z^2}.
 \label{eq:reader-beta-rational}
\end{equation}

\subsection{Threshold routing}

For $0<d<1-\sqrt3/2$, the exact high-demand threshold is
\begin{equation}
 z>e(d)
 \quad\Longrightarrow\quad
 G_{1-d}(z)\ge1-e(d).
 \label{eq:reader-high-threshold}
\end{equation}

\begin{lemma}[Threshold routing]
\label{lem:reader-threshold-routing}
Let $z_j=G_j(z_{j-1})$, where every $G_j$ is nondecreasing and extensive.
If $G_k=G_{1-d}$ and $z_{k-1}>e(d)$, then
\[
 z_j\ge1-e(d)
 \qquad(j\ge k).
\]
If a chain contains later occurrences of $G_{1-d_1}$ and $G_{1-d_2}$ and
\[
 z_0>e(d_1),
 \qquad
 \min\{e(d_1),e(d_2)\}<Q,
\]
then one of those two rows forces every subsequent output above $1-Q$.
\end{lemma}

\begin{proof}
Equation~\eqref{eq:reader-high-threshold} gives the first triggered output,
and extensivity preserves it.  For the two-threshold statement, if
$e(d_1)<Q$, trigger at the first row.  Otherwise
\[
 z_0>e(d_1)\ge Q>e(d_2),
\]
so trigger at the second row.
\end{proof}

The lemma is used only after Equation~\eqref{eq:reader-actual-propagation} has
identified each formal input as a lower bound for an actual row.

\subsection{One signed five-row chain}

Use the signed center data of
Proposition~\ref{prop:signed-center-normal-form}.  Put
\[
 X=R-\delta,
 \qquad
 H=\frac{\eta+\alpha+\delta}{2R},
 \qquad
 m=\min\left\{\frac\alpha R,\frac\delta W\right\}.
\]
The five nonsupercritical demands are
\[
 1-\frac\delta R,
 \quad1-\delta,
 \quad1-m,
 \quad1-\alpha,
 \quad1-\frac\alpha W.
\]
A gap in the normalized right trace starts the same actual chain
$1\to2\to3\to4\to5\to0$ for CE1 and CE2.  The common interface in
Proposition~\ref{prop:signed-one-gap-interface} reduces both classes to
\begin{equation}
 \left(G_{1-\delta}\circ G_{1-m}\circ G_{1-\alpha}\right)(H)
 >W+\delta.
 \label{eq:reader-signed-three-map}
\end{equation}
The only remaining distinction is the sign of
$\Delta_L=P-R\alpha-\delta$.

\begin{proposition}[Simplified CE1 one-gap obstruction]
\label{prop:reader-ce1-one-gap}
The CE1, $N_+=1$, all-$\Vdzero$, one-gap state is impossible.
\end{proposition}

\begin{proof}
Here $\Delta_L\le0$.  The low-root estimates give $e(\alpha)<X$ and
$H>\alpha$.  At the row with deficit $\alpha$, the Full labels are impossible.
The $L$ and $T_-$ labels immediately force an output above $1-X$.

On the selected $T_+$ label, the center inequalities force $X>1/2$.  Applying
\eqref{eq:reader-high-chord} first at deficit $\alpha$ and then the appropriate
high- or low-radial chord at deficit $m$ gives lower bounds
\[
 p_1>(2-4\alpha)H-(1-4\alpha)\alpha=:L_1,
\]
\[
 p_2>(2-5m)L_1-(1-5m)m=:L_2.
\]
The exact terminal calculation in Appendix~\ref{sec:cyclic-all-vd0} proves
\begin{equation}
 \delta<\frac1{10},
 \qquad
 p_2>L_2>2\delta+3\delta^2>e(\delta).
 \label{eq:reader-ce1-terminal}
\end{equation}
Threshold routing at the row with deficit $\delta$ therefore gives an output
above $1-X$, proving \eqref{eq:reader-signed-three-map}.  The common actual-row
interface then yields
\[
 a_0\ge Z>B_k(k/R)\ge a_0,
\]
a contradiction.  All inequalities used in the selected branch, including
the proof that $X>1/2$ and the terminal estimate
\eqref{eq:reader-ce1-terminal}, are proved in the cited appendix.
\end{proof}

\begin{proposition}[Simplified CE2 one-gap obstruction]
\label{prop:reader-ce2-one-gap}
The CE2, $N_+=1$, all-$\Vdzero$, exactly-one-gap state is impossible in either
orientation.
\end{proposition}

\begin{proof}
Here $\Delta_L>0$.  Proposition~\ref{prop:signed-ce2-one-gap} proves
\[
 R-\delta>e(\alpha),
 \qquad
 \min\{e(\alpha),e(\delta)\}<H.
\]
If $e(\alpha)<H$, threshold routing triggers at the row with deficit
$\alpha$.  Otherwise
\[
 R-\delta>e(\alpha)\ge H>e(\delta),
\]
and it triggers at the row with deficit $\delta$.  In either case the return
output exceeds $1-H$, contradicting the common supercritical-row bound.
Reflection exchanges $R$ with $W$, $\alpha$ with $\delta$, and reverses the
row order, so the second orientation is identical.
\end{proof}

\subsection{The remaining exact-demand branches}

The CE2 two-gap state uses the same endpoint deficits $\alpha/W$ and
$\delta/R$, but the two gaps share row $T_0$; the paired capped-loss theorem is
therefore retained.  The $N_+=0$ all-Vd0 proof uses the same one-side loss and
the common three-row boundary budget.  The T3-like high sheet is the universal
curve \eqref{eq:reader-psi}, with the rational parameter
\eqref{eq:reader-beta-rational}; only its branch-specific endpoint estimates
remain in Appendix~\ref{sec:t3-vd-demand-branches}.

For the one-Vd1/Vd2 branch, Lemma~\ref{lem:signed-small-slack} first gives
\[
 \alpha+\delta<\frac1{24},
 \qquad
 \alpha+\delta<\frac{\min\{R,W\}}6.
\]
These bounds prove the adjacent and nonadjacent placements in
Lemmas~\ref{lem:signed-vd-adjacent-placement} and
\ref{lem:signed-vd-nonadjacent-placement}.  The common hiding argument for a
T3-like or Vd1 adjacent rescuer is
Lemma~\ref{lem:signed-adjacent-rescuer}.

\begin{proposition}[Branches closed by exact demand]
\label{prop:reader-demand-branches}
Strategy~2 closes all routing-table rows assigned to exact demand propagation:
\begin{enumerate}
\item the $N_+=0$ all-$\Vdzero$ and T3-like-without-Vd1/Vd2 branches for
      $\CEone$ and $\CEtwo$;
\item every positive-gap $N_+=1$ all-$\Vdzero$ state;
\item the exactly-one-T3-like $N_+=1$ state with no Vd1/Vd2 role;
\item the complementary exact placements in the $\CEtwo$, $N_+=1$,
      exactly-one-Vd1/Vd2 branch.
\end{enumerate}
Singleton gaps are included.
\end{proposition}

\begin{proof}
The first item is Propositions~\ref{prop:nplus-zero-all-vd0} and
\ref{prop:nplus-zero-t3}.  For the second, the one-gap states are
Propositions~\ref{prop:reader-ce1-one-gap} and
\ref{prop:reader-ce2-one-gap}; the two-gap state and the no-gap transfer are
assembled in Proposition~\ref{prop:nplus-one-all-vd0}.  The third item is
Proposition~\ref{prop:nplus-one-one-t3}.  The fourth is
Proposition~\ref{prop:signed-ce2-one-vd-placements}.  In every positive-gap
chain, Equation~\eqref{eq:reader-actual-propagation} supplies the passage from
scalar maps to actual role reaches.
\end{proof}
